\documentclass[11pt]{ctexart}
\usepackage[a4paper,margin=1in]{geometry}
\usepackage{graphicx}
\usepackage{xcolor}
\usepackage{hyperref}
\definecolor{refgray}{gray}{0.45}
\title{\textbf{市场最新观点汇总}\\[0.4em]{\large 结构性分化与再定价成为市场主线，通胀压力从能源向核心领域扩散，风险偏好修复但空间有限。}}
\author{KC桌面 自动生成}
\date{260531}
\begin{document}
\maketitle
\section*{一页摘要}
\begin{itemize}
  \item 全球通胀格局正在经历结构性切换，核心通胀的隐性上调比能源价格冲击更值得警惕。
  \item 中国钢铁下游需求呈现结构性分化，建筑、机械、汽车等行业冷热不均，总量回暖掩盖了质量差异。
  \item 城市更新五年规划带来4-5万亿元增量投资，供给侧再定价逻辑正在重塑房地产投资框架。
  \item 美债长端反弹主要来自风险溢价压缩而非降息预期升温，后续空间有限。
  \item 全球石油需求下降约9\%，但替代的不可逆性意味着部分需求可能永久消失。
  \item AI基础设施融资并非瓶颈，资本中介结构的重塑才是关键挑战。
  \item 欧洲并购窗口不在大盘蓝筹，而在被低估的中小盘折价。
  \item 黄金的支撑力量——官方需求——正在衰减，利率敏感度回归。
  \item 印尼正在经历政策-资本-评级三重负反馈，卢比弱势可能持续。
  \item 日本国债曲线中段表现优于长端，市场焦点从通胀风险溢价转向实际经济数据。
\end{itemize}
\section{宏观与利率}
\textbf{全球利率市场的核心驱动力正从政策预期转向风险溢价压缩，但这一压缩空间已大幅收窄；通胀压力正从能源向核心领域扩散，央行政策路径面临重新定价。}

\begin{itemize}
  \item 美债长端收益率回落主要由期限溢价压缩驱动，市场对美联储利率路径的预期几乎没有变化，这意味着后续收益率进一步下行的空间有限。
  \item 全球核心通胀表面上持平于2.3\%，但多家经济体的核心通胀预测已被上调，中位数上调约0.5个百分点，美国与欧元区分别上调约0.4个百分点，涨价压力正在扩散。
  \item 全球PPI近期明显回升，从1.4\%升至2.8\%，生产者价格向消费者价格的传导可能尚未完全体现。
  \item 日本国债曲线中段表现优于长端，市场正在从对通胀风险溢价的过度定价转向关注实际消费者通胀数据，植田和男的讲话进一步为过热通胀担忧降温。
  \item 全球债券市场出现联动抛售，美、英、日十年期国债收益率单周分别上行21、24和23个基点，G7财长会议已将其列为正式议题。
  \item 欧洲利率仍有下行空间，欧央行6月加息概率高，叠加活动数据偏弱和商品价格缓解，欧元区利率走势相对扎实。
\end{itemize}
\begin{figure}[htbp]
\centering
\includegraphics[width=0.88\linewidth]{figures/fig_150.jpg}
\caption{Exhibit 2 - Exhibit 1: The recent US rally has been predominantly driven by a reset in Treasury term premium Change in 10y UST term premium/rate expectations, GS estimate}
\end{figure}
\begin{figure}[htbp]
\centering
\includegraphics[width=0.88\linewidth]{figures/fig_151.jpg}
\caption{Exhibit 2 - Exhibit 2: Implied vol is at the bottom of our range of estimates of fair Average Implied Vol vs fair value model range}
\end{figure}
\begin{figure}[htbp]
\centering
\includegraphics[width=0.88\linewidth]{figures/fig_152.jpg}
\caption{Exhibit 3 - Exhibit 3: Sovereign spread compression has run ahead of relief in rates levels and volatility Retracement from post-conflict peak / wides. Using GS fitted yields}
\end{figure}
\begin{figure}[htbp]
\centering
\includegraphics[width=0.88\linewidth]{figures/fig_014.jpg}
\caption{Figure 1 - \# Global Inflation (April)\textbackslash{}* Figure 1. Global Inflation (April)}
\end{figure}
\begin{figure}[htbp]
\centering
\includegraphics[width=0.88\linewidth]{figures/fig_015.jpg}
\caption{Figure 2 - © 2026 Citi Inc. No redistribution without Citi's written permission. Figure 2. Global Inflation (April)}
\end{figure}
\begin{figure}[htbp]
\centering
\includegraphics[width=0.88\linewidth]{figures/fig_016.jpg}
\caption{Figure 3 - © 2026 Citi Inc. No redistribution without Citi's written permission. Figure 3. Global Headline CPI Inflation: Country-Level Detail (April) 3-Month Change (Pct Pts)}
\end{figure}
\begin{figure}[htbp]
\centering
\includegraphics[width=0.88\linewidth]{figures/fig_017.jpg}
\caption{Figure 4 - © 2026 Citi Inc. No redistribution without Citi's written permission. Figure 4. Global Core CPI Inflation: Country-Level Detail (April) 3-Month Change (Pct Pts)}
\end{figure}
\begin{figure}[htbp]
\centering
\includegraphics[width=0.88\linewidth]{figures/fig_018.jpg}
\caption{Figure 5 - © 2026 Citi Inc. No redistribution without Citi's written permission. Figure 5. Global PPI Inflation: Country-Level Detail (March/April) © 2026 Citi Inc. No redistribution without Citi's written permission. Figure 6. DM Country-Level CPI Inflation (April)}
\end{figure}
\begin{figure}[htbp]
\centering
\includegraphics[width=0.88\linewidth]{figures/fig_202.jpg}
\caption{Exhibit 1 - Exhibit 1: Weekly change in yield in both JGB and}
\end{figure}
\begin{figure}[htbp]
\centering
\includegraphics[width=0.88\linewidth]{figures/fig_203.jpg}
\caption{Exhibit 2 - Exhibit 2: Japan Corporate Goods price Index}
\end{figure}
\begin{figure}[htbp]
\centering
\includegraphics[width=0.88\linewidth]{figures/fig_204.jpg}
\caption{Exhibit 3 - Exhibit 3: Nikkei CPI T Index (y/y) vs. Nationwide CPI food less fresh food}
\end{figure}
\begin{figure}[htbp]
\centering
\includegraphics[width=0.88\linewidth]{figures/fig_205.jpg}
\caption{Exhibit 4 - Exhibit 4: Nationwide CPI inflation excluding special factors}
\end{figure}
{\small\color{refgray}\textbf{References:} [R002] 全球通胀的真正转折点不在油价，而在核心通胀的隐性扩散; [R005] 市场低估了“风险溢价压缩”的边界，而非降息预期本身; [R010] 市场正在重新定价的不是通胀本身，而是通胀风险溢价; [R015] 黄金的真正问题不在利率，而在官方需求的结构性对冲能力正在衰减}

\section{AI与科技}
\textbf{AI基础设施投资周期正在重塑资本市场中介结构，融资并非瓶颈，但资本匹配效率是关键；半导体测试环节成为产业链最紧的瓶颈，存储市场呈现三层共存格局。}

\begin{itemize}
  \item 全球资本池约256万亿美元，10万亿美元的AI投资周期仅占约4\%，资本总量不是约束，关键在于如何将不同风险偏好、久期和流动性的资金匹配到AI基础设施的不同环节。
  \item AI基础设施融资正在催生多元化生态，保险资金偏好长期固收，养老金配置更广，私募股权擅长高风险项目，公共债券市场仍有17万亿美元的吸收能力。
  \item 半导体测试环节正在从成本中心转变为产能瓶颈，GPU测试时间从Blackwell的约850秒延长至Rubin的约1200秒，单颗芯片对测试设备的需求增加近40\%。
  \item 数据中心存储正在形成SSD+HDD+磁带的三层架构，企业级SSD出货容量同比暴增116.7\%，NL HDD出货容量同比增长14.3\%，两者都在向更高容量产品迁移。
  \item 中国三大运营商已正式推出AI Token计费方案，从按流量包月转向按Token用量计费，定价体系初步成型，AI变现节奏可能比市场预期快一到两个季度。
  \item 苹果正在规划将COB封装方案引入摄像头模组，先用于可穿戴设备，最早2028年可能上iPhone，这将重塑供应链竞争格局。
\end{itemize}
\begin{figure}[htbp]
\centering
\includegraphics[width=0.88\linewidth]{figures/fig_171.jpg}
\caption{Exhibit 1 - Exhibit 1: Deep and diverse capital pools to finance estimated \textbackslash{}\textasciitilde{}\textbackslash{}\$10T AI build-out this cycle ```mermaid graph TD}
\end{figure}
\begin{figure}[htbp]
\centering
\includegraphics[width=0.88\linewidth]{figures/fig_172.jpg}
\caption{Exhibit 2 - Exhibit 2: Every tech cycle delivers 10x more compute Computing Cycles Over Time, 1960s – Today Devices / Users (MM in Log Scale)}
\end{figure}
\begin{figure}[htbp]
\centering
\includegraphics[width=0.88\linewidth]{figures/fig_173.jpg}
\caption{Exhibit 3 - Exhibit 3: AI is following the path of past cycles Early Cycle Cloud vs. AI Era Capex (\$B)}
\end{figure}
\begin{figure}[htbp]
\centering
\includegraphics[width=0.88\linewidth]{figures/fig_174.jpg}
\caption{Exhibit 4 - Exhibit 4: Token usage is increasing non-linearly, as advancing capabilities and demand inflection points to a larger cycle than expected}
\end{figure}
\begin{figure}[htbp]
\centering
\includegraphics[width=0.88\linewidth]{figures/fig_175.jpg}
\caption{Exhibit 5 - Exhibit 5: Capex revisions reflect upside surprises and highlight the challenge of forecasting exponential growth}
\end{figure}
\begin{figure}[htbp]
\centering
\includegraphics[width=0.88\linewidth]{figures/fig_176.jpg}
\caption{Exhibit 6 - Exhibit 6: The Emerging AI Infrastructure Financing Stack \# The Emerging AI Infrastructure Financing Stack A multi-asset ecosystem matching capital to AI infrastructure risk, return and duration The AI build-out should be support}
\end{figure}
\begin{figure}[htbp]
\centering
\includegraphics[width=0.88\linewidth]{figures/fig_177.jpg}
\caption{Exhibit 7 - Exhibit 7: Global Individual and Institutional Capital Pools of \textbackslash{}\$256T represent \textbackslash{}\textasciitilde{}\textbackslash{}\$4\% of the estimated \textbackslash{}\$10T needed to finance AI build-out Asset Owner Pools (\$T)}
\end{figure}
\begin{figure}[htbp]
\centering
\includegraphics[width=0.88\linewidth]{figures/fig_182.jpg}
\caption{Exhibit 12 - Exhibit 12: We see \textbackslash{}\$3.2T of data center capex spend through 2028 Updated Data Center Capex – Growing Role of Credit Markets ```mermaid graph TD}
\end{figure}
\begin{figure}[htbp]
\centering
\includegraphics[width=0.88\linewidth]{figures/fig_208.jpg}
\caption{Exhibit 1 - Exhibit 1: HDD Production Volume YoY and Share Prices of TDK \& Nidec}
\end{figure}
\begin{figure}[htbp]
\centering
\includegraphics[width=0.88\linewidth]{figures/fig_209.jpg}
\caption{Exhibit 5 - Exhibit 5: HDD Supply Chain ```mermaid graph TD}
\end{figure}
\begin{figure}[htbp]
\centering
\includegraphics[width=0.88\linewidth]{figures/fig_258.jpg}
\caption{Exhibit 3 - Exhibit 3: Monthly new 5G base station additions: Apr 2026 deployment at 51k units, or \$16\textbackslash{}\%\$ YoY}
\end{figure}
\begin{figure}[htbp]
\centering
\includegraphics[width=0.88\linewidth]{figures/fig_259.jpg}
\caption{Exhibit 4 - Exhibit 4: 5G base station deployment by quarter: 1Q26 deployment decreased QoQ}
\end{figure}
\begin{figure}[htbp]
\centering
\includegraphics[width=0.88\linewidth]{figures/fig_260.jpg}
\caption{Exhibit 5 - Exhibit 5: Base station production was \$-19\textbackslash{}\%\$ YoY in Apr 2026 Exhibit 6: The number of newly installed 5G base stations in China was up to 51k units in Apr 2026 from 49k units in Mar 2026}
\end{figure}
{\small\color{refgray}\textbf{References:} [R008] 市场真正低估的不是AI融资缺口，而是资本中介结构的重塑; [R011] 数据中心存储的真正拐点，不在SSD对HDD的替代，而在供应链的再分配; [R012] 市场真正低估的不是需求，而是测试环节的结构性瓶颈; [R017] 三大运营商正式开卖AI“Token”：ARPU的下一轮增长，比市场预期的更早到来; [R020] 苹果拥抱COB封装，市场低估了供应链的“双赢”格局}

\section{能源与大宗商品}
\textbf{全球石油需求下降并非来自经济衰退，而是消费者主动选择替代品，部分需求可能永久消失；黄金的官方需求对冲能力正在衰减，利率敏感度回归。}

\begin{itemize}
  \item 中国石油需求可能下降约9\%（约150万桶/天），但公路运输指标并未明显走弱，消费者在价格信号下主动转向电动汽车和铁路出行。
  \item 中国高速公路电动汽车充电量在五一假期首日同比暴增55.6\%，1540万辆电动汽车上路占全部车辆的24\%，同比增加33\%。
  \item 欧洲可再生能源发电量大幅增长，即使油价接近120美元/桶，大部分国家电力价格反而跌入负值区间，化石能源空间被挤压。
  \item 黄金从4700美元附近持续回落，与原油走势背离，真正约束来自官方需求对冲能力的衰减——央行购金虽超预期，但ETF需求同比萎缩78\%。
  \item 2022年黄金曾经历270个基点的实际利率飙升，但被716吨的官方黄金需求摆动有效对冲，当前这一对冲机制正在发生结构性变化。
  \item 全球债券市场抛售正在重塑黄金的宏观定价框架，实际利率仍过低，不排除十年期美债收益率升至5\%的可能性。
\end{itemize}
\begin{figure}[htbp]
\centering
\includegraphics[width=0.88\linewidth]{figures/fig_163.jpg}
\caption{Figure 7 - Figure 1: Global energy stack}
\end{figure}
\begin{figure}[htbp]
\centering
\includegraphics[width=0.88\linewidth]{figures/fig_164.jpg}
\caption{Figure 3 - Figure 3: China electrified vehicle sales LHS: Number of electrified vehicles sold. RHS: share of total vehicle sales}
\end{figure}
\begin{figure}[htbp]
\centering
\includegraphics[width=0.88\linewidth]{figures/fig_165.jpg}
\caption{Figure 5 - Figure 5: China electric bus sales LHS: Number of buses sold. RHS: share of of total bus sales}
\end{figure}
\begin{figure}[htbp]
\centering
\includegraphics[width=0.88\linewidth]{figures/fig_166.jpg}
\caption{Figure 2 - Figure 2: Oil intensity of GDP Barrels of oil consumption per \textbackslash{}\$1000 of real GDP}
\end{figure}
\begin{figure}[htbp]
\centering
\includegraphics[width=0.88\linewidth]{figures/fig_167.jpg}
\caption{Figure 4 - Figure 4: China gasoline/diesel vehicle sales LHS: Number of gasoline/diesel vehicles sold. RHS: share of total vehicle sales}
\end{figure}
\begin{figure}[htbp]
\centering
\includegraphics[width=0.88\linewidth]{figures/fig_168.jpg}
\caption{Figure 6 - Figure 6: China electric truck sales LHS: Number of trucks sold. RHS: share of total truck sales}
\end{figure}
\begin{figure}[htbp]
\centering
\includegraphics[width=0.88\linewidth]{figures/fig_210.jpg}
\caption{Figure 1 - Figure 1: Gold diverging from oil}
\end{figure}
\begin{figure}[htbp]
\centering
\includegraphics[width=0.88\linewidth]{figures/fig_211.jpg}
\caption{Figure 2 - Figure 2: 10y US Treasury nominal yield more closely related to Dec'26 Brent than front month (where coefficient of determination is 0.51)}
\end{figure}
\begin{figure}[htbp]
\centering
\includegraphics[width=0.88\linewidth]{figures/fig_212.jpg}
\caption{Figure 2 - Figure 3: 10y real rates have followed Dec Fed pricing}
\end{figure}
\begin{figure}[htbp]
\centering
\includegraphics[width=0.88\linewidth]{figures/fig_213.jpg}
\caption{Figure 4 - Figure 4: Gold Dec'26 and Fed Dec'26}
\end{figure}
\begin{figure}[htbp]
\centering
\includegraphics[width=0.88\linewidth]{figures/fig_214.jpg}
\caption{Figure 5 - Figure 5: Gold demand weaker in key categories}
\end{figure}
\begin{figure}[htbp]
\centering
\includegraphics[width=0.88\linewidth]{figures/fig_215.jpg}
\caption{Figure 6 - Figure 6: Recycled gold in Q1 of 366 tonnes supports full year assumption}
\end{figure}
{\small\color{refgray}\textbf{References:} [R007] 市场真正低估的不是需求崩塌，而是替代的不可逆性; [R015] 黄金的真正问题不在利率，而在官方需求的结构性对冲能力正在衰减}

\section{地产与消费}
\textbf{城市更新规划带来供给侧再定价，增量投资规模相当于地产固投的66\%；钢铁下游需求结构性分化，建筑、机械、汽车冷热不均。}

\begin{itemize}
  \item 十五五城市更新规划隐含增量投资约4-5万亿元，相当于2026-2027年房地产固定投资的66\%，与上一轮棚改高峰期增量占比几乎持平。
  \item 城市更新重点方向包括危旧房改造翻倍至50万套、老旧街区/厂区改造增加67\%、地下管网改造增加约50\%，同时新增城市公园绿地改造等赛道。
  \item 5月钢铁下游PMI综合指数50.66\%，环比微升，但建筑行业新订单指数创近期高点52.3\%的同时生产指数却在下降，新增订单主要来自短周期、快回款项目。
  \item 机械行业PMI连续第二个月保持51\%以上，机床和变压器市场订单火爆，但农机因春耕结束需求走弱。
  \item 汽车行业PMI环比跳升1.6个百分点至50.84\%，五一假期和促销拉动新订单回暖，新能源车继续跑赢燃油车。
  \item 消费信贷数据与情绪指标之间出现裂痕，密歇根消费者信心指数创纪录新低，但实际违约率尚未明显恶化。
\end{itemize}
\begin{figure}[htbp]
\centering
\includegraphics[width=0.88\linewidth]{figures/fig_001.jpg}
\caption{Figure 1 - Cynthia Wu Figure 1. Steel PMI Composite Index © 2026 Citi Inc. No redistribution without Citi's written permission. \# Construction Industry Figure 2. Construction Industry PMI}
\end{figure}
\begin{figure}[htbp]
\centering
\includegraphics[width=0.88\linewidth]{figures/fig_002.jpg}
\caption{Figure 3 - © 2026 Citi Inc. No redistribution without Citi's written permission. Construction PMI was 50.7\%, up 0.1\% MoM and 0.9ppt YoY. On a seasonally adjusted basis, the reading was 49.5\%, up 0.6ppt MoM and 0.8ppt YoY. The production index was 50.1\%, down 1.2ppt MoM a}
\end{figure}
\begin{figure}[htbp]
\centering
\includegraphics[width=0.88\linewidth]{figures/fig_003.jpg}
\caption{Figure 4 - © 2026 Citi Inc. No redistribution without Citi's written permission. Figure 4. Composite index-Construction}
\end{figure}
\begin{figure}[htbp]
\centering
\includegraphics[width=0.88\linewidth]{figures/fig_004.jpg}
\caption{Figure 6 - © 2026 Citi Inc. No redistribution without Citi's written permission. Machinery PMI was 51\%, up 0.3\% MoM and +1.7ppt YoY. On a seasonally adjusted basis, the reading was 50.3\%, down 0.2ppt MoM and up 1.6ppt YoY. Production index was 51.6\%, up 0.9ppt MoM and 2.}
\end{figure}
\begin{figure}[htbp]
\centering
\includegraphics[width=0.88\linewidth]{figures/fig_005.jpg}
\caption{Figure 7 - © 2026 Citi Inc. No redistribution without Citi's written permission. Figure 7. Composite index - Machinery}
\end{figure}
\begin{figure}[htbp]
\centering
\includegraphics[width=0.88\linewidth]{figures/fig_006.jpg}
\caption{Figure 9 - © 2026 Citi Inc. No redistribution without Citi's written permission. Auto PMI was 50.8\%, up 1.6\% MoM and +0.2ppt YoY. On a seasonally adjusted basis, the reading was 50\%, up 0.4ppt MoM and flattish YoY. Production index was 51.4\%, up 2.1ppt MoM and up 0.5ppt }
\end{figure}
\begin{figure}[htbp]
\centering
\includegraphics[width=0.88\linewidth]{figures/fig_007.jpg}
\caption{Figure 10 - © 2026 Citi Inc. No redistribution without Citi's written permission. Figure 10. Composite index - Automobile}
\end{figure}
\begin{figure}[htbp]
\centering
\includegraphics[width=0.88\linewidth]{figures/fig_025.jpg}
\caption{Exhibit 1 - Exhibit 1: We see acceleration of renovation for urban dilapidated housing, old neighborhoods and factory areas, historical building and urban underground pipeline network, with several niche fields of urban renewal newly introduce}
\end{figure}
\begin{figure}[htbp]
\centering
\includegraphics[width=0.88\linewidth]{figures/fig_026.jpg}
\caption{Exhibit 2 - Comparison of property FAI and incremental FAI from shantytown redevelopment and urban renewal}
\end{figure}
\begin{figure}[htbp]
\centering
\includegraphics[width=0.88\linewidth]{figures/fig_283.jpg}
\caption{Exhibit 13 - Exhibit 13: The University of Michigan Consumer Sentiment Index fell to a new record low in May, with an average of 48\% of consumers citing high prices as eroding their financial wellbeing University of Michigan Consumer sentiment i}
\end{figure}
\begin{figure}[htbp]
\centering
\includegraphics[width=0.88\linewidth]{figures/fig_284.jpg}
\caption{Exhibit 14 - Exhibit 14: Consumer sentiment is worse for younger and lower income consumers, but only marginally so Consumer sentiment by income tercile (left panel) and age bracket (right panel)}
\end{figure}
\begin{figure}[htbp]
\centering
\includegraphics[width=0.88\linewidth]{figures/fig_285.jpg}
\caption{Exhibit 15 - Exhibit 15: Conference Board data came in at the 40th percentile since 1980, far from the record lows of the University of Michigan survey}
\end{figure}
\begin{figure}[htbp]
\centering
\includegraphics[width=0.88\linewidth]{figures/fig_286.jpg}
\caption{Exhibit 16 - Exhibit 16: Higher income consumers have experienced a much greater improvement in their financial situation vs. lower income consumers Percentage of Respondents Concerned About Making Ends Meet in 0–6 Months (left panel) and 7–12 m}
\end{figure}
\begin{figure}[htbp]
\centering
\includegraphics[width=0.88\linewidth]{figures/fig_287.jpg}
\caption{Exhibit 17 - Exhibit 17: With tax refund season fully in the rearview, seasonal drops in delinquencies have largely been in line with historical norms Drops in D30+ rates during tax season over time segmented by sector}
\end{figure}
\begin{figure}[htbp]
\centering
\includegraphics[width=0.88\linewidth]{figures/fig_288.jpg}
\caption{Exhibit 18 - Exhibit 18: Despite decreases in sentiment, delinquencies are not yet showing any clear impact from the conflict in the middle east D30+ rates for a sample of 142 subprime auto deals originated prior to 2024}
\end{figure}
{\small\color{refgray}\textbf{References:} [R001] 钢铁下游的真实信号：不是全面复苏，而是结构性分化正在加速; [R003] 城市更新的真实规模：市场低估的不是需求，而是供给侧的再定价; [R018] 市场真正低估的不是风险，而是结构性的相对价值错位}

\section{信用与风险偏好}
\textbf{高收益债市场对评级迁移风险的定价系统性低估了意外概率，美欧市场走向截然不同的净迁移方向；欧洲并购窗口不在大盘而在中小盘折价。}

\begin{itemize}
  \item 自2010年以来，美欧市场约一半的坠落天使在事件发生前未被评级机构明确标记为负面展望或降级观察，市场对评级迁移风险的定价可能系统性低估了意外概率。
  \item 美国高收益市场2026年上升之星预测上调至750亿美元，坠落天使预测为600亿美元，市场将连续多年处于净流入状态。
  \item 欧洲市场则相反，上升之星预测为200亿欧元，坠落天使预测为300亿欧元，若兑现将是2020年以来首次坠落天使规模超过上升之星。
  \item 欧洲并购活动低迷，占全球份额降至不足三分之一，但资产剥离意愿创历史纪录，企业主动释放优质业务单元。
  \item 欧洲中小盘股相对大盘的估值折价接近历史最深水平，STOXX Small相对Large的12个月远期市盈率折价显著。
  \item 私募退出压力加大，IPO窗口仍窄，手握充足现金的买家有机会以合理估值收购优质资产。
\end{itemize}
\begin{figure}[htbp]
\centering
\includegraphics[width=0.88\linewidth]{figures/fig_027.jpg}
\caption{Exhibit 1 - Exhibit 1: On a spread basis, USD BB-rated bonds have meaningfully outperformed B-rated bonds since late 2025 OAS ratio of the B and BB subsets of the Bloomberg USD HY Corporate indices}
\end{figure}
\begin{figure}[htbp]
\centering
\includegraphics[width=0.88\linewidth]{figures/fig_028.jpg}
\caption{Exhibit 2 - Exhibit 2: In contrast, USD BB vs. BBB spreads have been rangebound OAS ratio of the BB and BBB subsets of the Bloomberg USD HY and IG Corporate indices}
\end{figure}
\begin{figure}[htbp]
\centering
\includegraphics[width=0.88\linewidth]{figures/fig_029.jpg}
\caption{Exhibit 3 - Exhibit 3: In the EUR market, BB spreads have also outperformed the B-rated cohort OAS ratio of the B and BB subsets of the Bloomberg Pan Euro HY Corporate indices}
\end{figure}
\begin{figure}[htbp]
\centering
\includegraphics[width=0.88\linewidth]{figures/fig_030.jpg}
\caption{Exhibit 4 - Exhibit 4: EUR BB vs. BBB spreads remain near the mid-to-low end of the historical range OAS ratio of the BB and BBB subsets of the Bloomberg Pan Euro HY and IG Corporate indices}
\end{figure}
\begin{figure}[htbp]
\centering
\includegraphics[width=0.88\linewidth]{figures/fig_031.jpg}
\caption{Exhibit 5 - Exhibit 5: A high share of fallen angels in the EUR and USD markets were not rated BBB or BBB- on Negative outlook or Downgrade Watch prior to the event Share of bonds in each rating bucket that became fallen angels within 12 month}
\end{figure}
\begin{figure}[htbp]
\centering
\includegraphics[width=0.88\linewidth]{figures/fig_032.jpg}
\caption{Exhibit 6 - Exhibit 6: Similarly, a high share of rising stars in the EUR and USD markets were not BB+ or BB on Positive Outlook or Upgrade Watch prior to the event Share of bonds in each rating bucket that became rising stars within 12 months}
\end{figure}
\begin{figure}[htbp]
\centering
\includegraphics[width=0.88\linewidth]{figures/fig_033.jpg}
\caption{Exhibit 8 - Exhibit 7: In the EUR market, fallen angels are set to outpace rising stars for the first time since 2020. Fallen angel and rising star notionals in the EUR market (left panel) and USD market (right panel) EUR Notional for Fallen}
\end{figure}
\begin{figure}[htbp]
\centering
\includegraphics[width=0.88\linewidth]{figures/fig_034.jpg}
\caption{Exhibit 8 - Exhibit 8: Neither market is pricing significant fallen angel risk}
\end{figure}
\begin{figure}[htbp]
\centering
\includegraphics[width=0.88\linewidth]{figures/fig_216.jpg}
\caption{Exhibit 7 - Exhibit 1: Mega-deals dominate the M\&A recovery 12m rolling sum of M\&A deals globally}
\end{figure}
\begin{figure}[htbp]
\centering
\includegraphics[width=0.88\linewidth]{figures/fig_217.jpg}
\caption{Exhibit 2 - Exhibit 2: European M\&A remains subdued 12m rolling sum of M\&A deals in Europe}
\end{figure}
\begin{figure}[htbp]
\centering
\includegraphics[width=0.88\linewidth]{figures/fig_218.jpg}
\caption{Exhibit 3 - Exhibit 3: Asset disposals rise to a record strategic priority UK CFO Survey: Share of respondents prioritising asset disposals over the next 12 months (\%)}
\end{figure}
\begin{figure}[htbp]
\centering
\includegraphics[width=0.88\linewidth]{figures/fig_219.jpg}
\caption{Exhibit 4 - Exhibit 4: Activity is increasingly driven by cross-border flows 12m rolling sum of M\&A deals in Europe}
\end{figure}
\begin{figure}[htbp]
\centering
\includegraphics[width=0.88\linewidth]{figures/fig_220.jpg}
\caption{Exhibit 5 - Exhibit 5: Small caps tend to be correlated with our European M\&A Candidates basket Price performance of GS M\&A Candidates Basket (GSTRACQN) vs. STOXX 600 and Small vs. Large caps (y/y chg, \%).}
\end{figure}
\begin{figure}[htbp]
\centering
\includegraphics[width=0.88\linewidth]{figures/fig_221.jpg}
\caption{Exhibit 6 - Exhibit 6: The M\&A Candidates basket (GSTRACQN) trades at a heavy discount to the market GSTRACQN vs. SXXP}
\end{figure}
\begin{figure}[htbp]
\centering
\includegraphics[width=0.88\linewidth]{figures/fig_222.jpg}
\caption{Exhibit 7 - Exhibit 7: STOXX Small are close to the deepest discount historically against STOXX Large 12m fwd P/E from Small caps vs. Large caps (premium/discount)}
\end{figure}
{\small\color{refgray}\textbf{References:} [R004] 高收益债市场真正低估的不是违约风险，而是评级迁移的“意外率”; [R016] 欧洲并购的真正窗口不在大盘，而在被市场遗忘的小盘折价}

\section{区域市场与外汇}
\textbf{印尼正在经历政策-资本-评级三重负反馈，卢比弱势可能持续；美加贸易谈判的结构性不对等对加元构成压力；人民币仍被低估。}

\begin{itemize}
  \item 印尼卢比的弱势正从地缘冲击下的被动贬值演变为国内政策选择驱动的主动走弱，政府推动自然资源出口集中化、更严格的出口收益留存规定打击了外资信心。
  \item MSCI和FTSE Russell因股权集中度问题剔除印尼股票，预计导致约17亿美元被动资金外流，即将到来的市场准入审查可能引发额外约40亿美元潜在流出。
  \item 做多新加坡元兑印尼卢比的交易信心被上调至最高等级，新加坡GDP增速上修至6\%，央行7月可能再次收紧货币政策。
  \item 美加贸易谈判的分歧远大于美墨，2018-2019年贸易冲突新闻占比每上升10个百分点，美元/加元月度变化约上升1.5\%，加元是北美贸易政策不确定性的第一接收器。
  \item 贸易冲击通过加拿大央行的反应函数传导至加元，2018年加拿大央行明确表示利率可能需要维持在低于中性水平，直到贸易不确定性消散。
  \item 人民币仍被低估约10\%，基于四种估值模型，华为芯片技术突破和中美关系缓和预期可能吸引更多资本流入。
\end{itemize}
\begin{figure}[htbp]
\centering
\includegraphics[width=0.88\linewidth]{figures/fig_156.jpg}
\caption{Exhibit 1 - Exhibit 1: News coverage suggests that negotiations with Canada were more contentious than with Mexico}
\end{figure}
\begin{figure}[htbp]
\centering
\includegraphics[width=0.88\linewidth]{figures/fig_157.jpg}
\caption{Exhibit 2 - We use Bloomberg News Trends to produce a count of news stories containing “trade war”, Canada (or Mexico), and USA, and separately “trade”, Canada (or Mexico), and USA. We then take the “trade war” share of total trade headlines for each country. We have foun}
\end{figure}
\begin{figure}[htbp]
\centering
\includegraphics[width=0.88\linewidth]{figures/fig_158.jpg}
\caption{Exhibit 3 - Exhibit 3: ...but had a more limited impact on MXN}
\end{figure}
\begin{figure}[htbp]
\centering
\includegraphics[width=0.88\linewidth]{figures/fig_159.jpg}
\caption{Exhibit 4 - Exhibit 4: We find markets price limited trade-specific premium beyond standard cyclical drivers in our GSBEER models}
\end{figure}
\begin{figure}[htbp]
\centering
\includegraphics[width=0.88\linewidth]{figures/fig_160.jpg}
\caption{Exhibit 5 - Exhibit 5: During the 2025 fentanyl-related tariffs and Liberation Day uncertainty, the moves in USD/CAD were largely explained by rate differentials}
\end{figure}
\begin{figure}[htbp]
\centering
\includegraphics[width=0.88\linewidth]{figures/fig_161.jpg}
\caption{Exhibit 6 - Exhibit 6: Trade headlines have tended to move USD-CAD rate differentials and USD/CAD together}
\end{figure}
\begin{figure}[htbp]
\centering
\includegraphics[width=0.88\linewidth]{figures/fig_162.jpg}
\caption{Exhibit 7 - Exhibit 7: Rising trade tensions are also reflected in low-delta options, consistent with stronger demand for protection against policy-driven tail risks}
\end{figure}
{\small\color{refgray}\textbf{References:} [R006] 市场真正低估的不是关税，而是美加贸易谈判的结构性不对等; [R009] 市场误读了监管意图：收紧的不是跨境资金，而是非法账户通道; [R014] 市场真正低估的不是地缘风险，而是印尼结构性抛售的持续性}

\section*{结语}
综合来看，当前市场主线已从总量复苏转向结构性分化与再定价。通胀压力从能源向核心领域扩散、AI投资重塑资本中介结构、能源替代的不可逆性、以及区域市场的自我强化循环，均提示投资者需要超越短期叙事，关注更深层的结构性变化。风险偏好修复虽在推进，但空间有限，且不同资产和区域之间的分化正在加剧。
\vfill
{\small\color{refgray}Personal reading notes and learning share only. Not investment advice.}
\end{document}
